\documentclass[11pt]{amsart}
\usepackage{amsmath,amssymb,amsthm}

\newtheorem{theorem}{Theorem}[section]
\newtheorem{lemma}[theorem]{Lemma}
\newtheorem{proposition}[theorem]{Proposition}
\newtheorem{corollary}[theorem]{Corollary}
\newtheorem{conjecture}[theorem]{Conjecture}
\theoremstyle{definition}
\newtheorem{definition}[theorem]{Definition}
\theoremstyle{remark}
\newtheorem*{remark}{Remark}

\title{A sum--reciprocal inequality and two applications}
\author{L. V. Instance}
\address{School of Notional Analysis, Imaginary University (fictional)}
\email{l.v.instance@example.org}
\date{}

\begin{document}

\begin{abstract}
We give a short derivation of the inequality
$(x_1+\dots+x_n)(1/x_1+\dots+1/x_n)\ge n^2$ for positive reals from the
Cauchy--Schwarz inequality, deduce the comparison of the arithmetic and
harmonic means and Nesbitt's inequality, and record an elementary bound on
$(1+1/n)^n$ obtained from Bernoulli's inequality.
\end{abstract}

\maketitle

\section{Introduction}

The inequality in Theorem~\ref{thm:main} below is a standard first
application of the Cauchy--Schwarz inequality~\cite{Cau}. This note is
intended as a self-contained reference for it and for a few consequences we
use elsewhere. Section~\ref{sec:sr} contains the main theorem and its
corollaries, including Nesbitt's inequality~\cite{Nes}; Section~\ref{sec:bern}
treats the sequence $(1+1/n)^n$.

\section{The sum--reciprocal inequality}\label{sec:sr}

Throughout this section $n\ge 1$ is an integer and $x=(x_1,\dots,x_n)$ is an
$n$-tuple of real numbers; all sums run over $i=1,\dots,n$.

\begin{definition}\label{def:means}
Let $x$ be an $n$-tuple of positive reals. Its \emph{arithmetic mean} and its
\emph{harmonic mean} are
\[
  A(x)=\frac1n\sum x_i,\qquad H(x)=\frac1n\sum\frac{1}{x_i}.
\]
\end{definition}

\begin{lemma}\label{lem:2ab}
For all real $a,b$ we have $a^2+b^2\ge 2ab$.
\end{lemma}

\begin{proof}
$a^2+b^2-2ab=(a-b)^2\ge 0$.
\end{proof}

\begin{lemma}\label{lem:recip}
For every real $x\ne 0$ we have $x+\frac1x\ge 2$.
\end{lemma}

\begin{proof}
Multiplying by $x$, the claim becomes $x^2+1\ge 2x$, which is
Lemma~\ref{lem:2ab} with $a=x$, $b=1$.
\end{proof}

\begin{lemma}[Cauchy--Schwarz~\cite{Cau}]\label{lem:CS}
For real $n$-tuples $a,b$ we have
$\bigl(\sum a_ib_i\bigr)^2\le\bigl(\sum a_i^2\bigr)\bigl(\sum b_i^2\bigr)$.
\end{lemma}

\begin{proof}
If all $a_i=0$ both sides vanish. Otherwise the quadratic polynomial
$q(t)=\sum(a_it+b_i)^2=t^2\sum a_i^2+2t\sum a_ib_i+\sum b_i^2$ is non-negative
for all real $t$ and has positive leading coefficient, so its discriminant
$4(\sum a_ib_i)^2-4\sum a_i^2\sum b_i^2$ is at most $0$.
\end{proof}

\begin{lemma}\label{lem:qm}
Let $x_1,\dots,x_n$ be real numbers with $x_i\ge 0$ for all $i$. Then
$\bigl(\sum x_i\bigr)^2\le n\sum x_i^2$.
\end{lemma}

\begin{proof}
Apply Lemma~\ref{lem:CS} with $a_i=x_i$ and $b_i=x_i$ for all $i$, and use
$\sum 1^2=n$.
\end{proof}

\begin{theorem}[Main theorem]\label{thm:main}
Let $n\ge 1$ and let $x_1,\dots,x_n$ be positive real numbers. Then
\[
  \Bigl(\sum x_i\Bigr)\Bigl(\sum\frac{1}{x_i}\Bigr)\ge n^2 .
\]
\end{theorem}

\begin{proof}
Apply Lemma~\ref{lem:CS} with $a_i=\sqrt{x_i}$ and $b_i=1/\sqrt{x_i}$. Then
$a_ib_i=1$, so the left-hand side of Lemma~\ref{lem:CS} is $n^2$, while the
right-hand side is $\bigl(\sum x_i\bigr)\bigl(\sum 1/x_i\bigr)$.
\end{proof}

\begin{corollary}[AM--HM]\label{cor:amhm}
For every $n$-tuple $x$ of positive reals, $A(x)\ge H(x)$.
\end{corollary}

\begin{proof}
Dividing Theorem~\ref{thm:main} by $n\sum\frac{1}{x_i}$ gives
$\frac1n\sum x_i\ge n\big/\sum\frac{1}{x_i}$, which is the claim.
\end{proof}

\begin{corollary}\label{cor:vac}
Let $x_1,\dots,x_n$ be positive reals with $\sum x_i=1$ and
$\sum 1/x_i<n^2$. Then the $x_i$ are not all equal.
\end{corollary}

\begin{proof}
If $x_i=1/n$ for all $i$, then $\sum 1/x_i=n^2$, contrary to the hypothesis.
\end{proof}

\begin{proposition}[Nesbitt~\cite{Nes}]\label{prop:nesbitt}
For all positive reals $a,b,c$,
\[
  \frac{a}{b+c}+\frac{b}{c+a}+\frac{c}{a+b}\ge\frac32 .
\]
\end{proposition}

\begin{proof}
Put $s=a+b+c$ and apply Theorem~\ref{thm:main} with $n=3$ and
$(x_1,x_2,x_3)=(b+c,\,c+a,\,a+b)$, whose sum is $2s$. This gives
$\sum\frac{1}{b+c}\ge\frac{9}{2s}$. Since
$\frac{a}{b+c}=\frac{s}{b+c}-1$, the left-hand side equals
$s\sum\frac{1}{b+c}-3\ge\frac92-3=\frac32$.
\end{proof}

\section{Bernoulli's inequality and $(1+1/n)^n$}\label{sec:bern}

We use Bernoulli's inequality in the form found in~\cite{HLP}: \emph{for every
real $x$ and every $n\in\mathbb{N}$,}
\[
  (1+x)^n\ge 1+nx .
\]
For $n=0$ both sides equal $1$, and the general case follows by induction on
$n$.

\begin{proposition}\label{prop:e}
For every $n\in\mathbb{N}$ we have $\bigl(1+\frac1n\bigr)^n\ge 2$.
\end{proposition}

\begin{proof}
Apply Bernoulli's inequality with $x=1/n$:
$\bigl(1+\frac1n\bigr)^n\ge 1+n\cdot\frac1n=2$.
\end{proof}

\begin{remark}
Proposition~\ref{prop:e} together with Lemma~\ref{lem:recip} shows that
$(1+\frac1n)^n+(1+\frac1n)^{-n}\ge 2$ as well, but this weaker statement
already follows from Lemma~\ref{lem:recip} alone.
\end{remark}

\begin{thebibliography}{9}
\bibitem{Cau} A.-L.~Cauchy, \emph{Cours d'analyse de l'\'Ecole Royale
Polytechnique}, Paris, 1821.
\bibitem{HLP} G.~H.~Hardy, J.~E.~Littlewood and G.~P\'olya,
\emph{Inequalities}, Cambridge University Press, 1934.
\bibitem{Nes} A.~M.~Nesbitt, Problem 15114, \emph{Educational Times}
\textbf{3} (1903), 37--38.
\end{thebibliography}

\end{document}
